\subsection{Interpreting Cross-Census Disruption}

The 66\% mean county disruption rate reflects \emph{demographic responsiveness}, not algorithmic instability. When 21--42\% of census tracts are split or merged between 2010 and 2020, METIS re-optimises from a substantially different graph. The algorithm faithfully incorporates updated population data; boundaries that were optimal for 2010 demographics are not necessarily optimal for 2020 demographics.

\textbf{Georgia versus Mississippi:} The 50-point difference in top-level split stability (86\% vs.\ 35\%) illustrates state-specific effects. Georgia's top-level north-south divide (Metro Atlanta vs.\ rest of state) is geographically stable even as population shifted toward the suburbs. Mississippi's split is more demographically sensitive: the Delta region's population declined relative to the Jackson metro, causing METIS to redraw the primary bisection axis.

\textbf{N-way comparison:} The original paper design included a recursive vs.\ n-way comparison. N-way run outputs for 2010 are not available in the current pipeline; this comparison is deferred to a Phase 2 run.

\subsection{Algorithmic Properties vs.\ Administrative Continuity}

Practitioners and courts distinguish between two types of stability:

\begin{enumerate}
    \item \textbf{Algorithmic consistency}: Does the algorithm apply the same objective function across census cycles? \textbf{Yes} — bisect uses identical METIS parameters (geographic edge weights, $\pm$0.5\% balance tolerance) in both 2010 and 2020 runs.

    \item \textbf{Geographic continuity}: Do district boundaries change between cycles? \textbf{Substantially} — 66\% county disruption reflects re-optimisation for new demographics.
\end{enumerate}

For litigation support, algorithmic consistency is the stronger property. A court cannot accuse bisect of partisan manipulation if the same algorithm with the same parameters produces both the 2010 and 2020 maps; any boundary changes are attributable to demographic change, not strategic intervention.

\textbf{Hybrid Approach:} One could seed 2030 redistricting with 2020's recursive tree structure, adjusting only lower-level splits. This would preserve algorithmic consistency while reducing unnecessary boundary disruption beyond what demographics require.

\subsection{Implications for 2030 Redistricting}

Our results generate a revised prediction for the 2030 census cycle:

\textbf{Prediction 1:} States with large demographic shifts (fast-growing metros, declining rural areas) will show higher cross-census disruption than states with stable demographic distributions regardless of algorithm choice.

\textbf{Prediction 2:} Preserving 2020's top-level split in 2030 will yield higher stability than re-running recursive bisection from scratch.

\textbf{Prediction 3:} The stability advantage will correlate with demographic volatility—states with rapid demographic change (e.g., Georgia, North Carolina) will benefit more from recursive methods than stable states (e.g., Vermont).

When 2030 census data arrives, researchers can test these predictions directly, validating or refuting our mechanistic explanation.

\subsection{Generalization Beyond Redistricting}

The performance-stability tradeoff extends to any domain requiring periodic re-partitioning:

\textbf{School Districts:} As enrollment shifts, school boundaries must adjust every 5-10 years. Hierarchical methods could minimize student reassignment.

\textbf{Service Territories:} Hospitals, fire departments, and police precincts rebalance catchments as populations grow. Stability reduces confusion and maintains institutional knowledge.

\textbf{Sales Regions:} Companies periodically rebalance sales territories. Hierarchical methods preserve regional organization while adjusting local boundaries.

\textbf{General Principle:} When re-partitioning frequency is low relative to inter-partition timescale, stability becomes valuable. Hierarchical methods provide continuity through structural preservation, while flat methods optimize each instance independently at the cost of disruption.

\subsection{Limitations of Current Analysis}

\textbf{Two Cycles Only:} We analyze 2010-2020 but lack 2000 data for three-cycle validation. Adding 2000 would strengthen claims about long-term stability.

\textbf{Algorithm-Specific:} Results apply to METIS-based partitioning. Other algorithms (spectral clustering, simulated annealing) may exhibit different stability properties.

\textbf{Tract Boundary Changes:} Census Bureau relationship files imperfectly map 2010-2020 tract changes. Some splits/merges introduce noise in stability metrics.

\textbf{Five States:} Limited to southern states with VRA scrutiny. Generalization to all 50 states requires additional experiments.
